\subsection{A lower bound after preprocessing the keys and values}
\label{sec:indexedlower}

The context-probe bounds proved earlier concern preprocessing that is
independent of the context.  An index built during prefill is a different
model.  The index can already contain every input sign, so changing one
sign can change the index before a query begins.  The earlier coupling
argument gives no query lower bound in this model.

We therefore use randomized OVH (Section~\ref{sec:index-main}) in
the following precise form: for each $\eta>0$ there is a constant $c>0$
such that no randomized algorithm, with error at most $1/3$, decides
whether two sets of $N$ vectors in $\{0,1\}^{\lceil c\log_2N\rceil}$
contain an orthogonal pair in $O(N^{2-\eta})$ bit operations.
The randomized form is needed because sampling a bounded-logit token
gives a randomized decision algorithm.  The standard split-and-list
reduction from randomized SETH gives randomized OVH.  The constant-gap
attention obstruction with polynomial prefill is also established by
Haverbeck et al.~\cite[Theorem 4.1 and Lemma A.5]{haverbeck2026attention}; the
direct reduction below keeps sign-valued coordinates and the
finite-bit stopping conditions explicit.

Our indexed attention problem has a static cache
$\mathcal C=((k_i,v_i))_{i=0}^{T-1}$ and an online query $q$.  It must
sample the binary token with probability
\begin{equation}
 \pi_{\mathcal C}(q)
 =\frac{1+\tanh(\tanh u_{\mathcal C}(q))}{2},
 \qquad
 u_{\mathcal C}(q)
 =\frac{\sum_{i=0}^{T-1}e^{\langle q,k_i\rangle/\sqrt n}v_i}
 {\sum_{i=0}^{T-1}e^{\langle q,k_i\rangle/\sqrt n}}.
 \label{eq:indexedtarget}
\end{equation}
This is a scalar attention coordinate, a unit tanh row, and a softmax
over the two logits $\pm\tanh u_{\mathcal C}(q)$.  Both logits lie in
$[-1,1]$ whenever $|v_i|\le1$.

\begin{theorem}[Polynomial prefill does not give uniformly sublinear queries]
\label{thm:indexedlower}
Assume randomized OVH.  For every fixed $a\ge1$, $b\ge0$, and
$\epsilon\in(0,1)$, there is a constant $C>0$ such that no uniform
randomized data structure has both of the following guarantees on all
caches of width $n\le C\log^4(T+2)$:
\begin{enumerate}
 \item it preprocesses the complete cache in expected
 $O(T^a(n+\log T)^b)$ bit operations;
 \item for every query it samples \eqref{eq:indexedtarget} in expected
 $O(T^{1-\epsilon}(n+\log T)^b)$ additional bit operations.
\end{enumerate}
The conclusion already holds when $n$ is a perfect square,
$q\in\{0,1\}^n$, $k_i\in\{-1,0\}^n$, and $v_i\in\{-1,1\}$.
In particular $\|q\|_2,\|k_i\|_2\le\sqrt n$ and all input numbers have
constant bit length.  The only score scaling is the standard
$1/\sqrt n$.

The same conclusion holds if the marginal law of each sample may have
total variation distance at most $1/64$ from the required token law.
It suffices to require the query guarantee for a fresh query to a fresh
index.  Correctness conditional on a random prefill state, or on an
adaptive query history, is not needed for this lower bound.
\end{theorem}

The constant $C$ is allowed to depend on the proposed preprocessing
degree and the claimed query saving.  This is the usual uniform
quantification in indexed fine-grained lower bounds.  The theorem
excludes a uniform $\operatorname{poly}(n,\log T)$ query algorithm after
any fixed polynomial prefill budget.  It does not give an unconditional
query lower bound, nor a lower bound for each individual cache.

\begin{remark}[Translating entry bounds into width]
The parameter regimes can be compared directly. Suppose old vectors
have dimension $d$, coordinate magnitudes at most $B_0$, and scores
$\langle q,k\rangle/d$. Choose an integer $H\ge\max\{1,B_0^2/d\}$.
Divide each coordinate by $\sqrt{dH}$ and repeat it $dH^2$ times.
The new width is $n=(dH)^2$, the coordinate bound is one, and
\[
 \frac{\langle q_{\rm new},k_{\rm new}\rangle}{\sqrt n}
 =\frac{dH^2}{dH\sqrt n}\langle q,k\rangle
 =\frac{\langle q,k\rangle}{d}.
\]
Thus $d=O(\log T)$ and $B_0=O(\sqrt{\log T})$, with
$B_0^2/d=O(1)$, translate to $n=O(\log^2T)$ at standard scaling.
Rounding the new coordinates within
$\eta\le1/(100\sqrt n)$, while keeping their magnitudes at most one,
changes every score by at most
$\sqrt n(2\eta+\eta^2)\le0.03$.
The dummy keys can remain exactly zero. These bounded perturbations
preserve a scalar separation tending to zero and one, since they
multiply all non-dummy exponential weights by factors in
$[e^{-0.03},e^{0.03}]$. Replacing scalar values $v\in\{0,1\}$ by
$2v-1$ then gives attention mean $u=2y-1$. Once the two cases have
$u\ge1/2$ and $u\le-1/2$, the token map in
\eqref{eq:indexedtarget} separates them by $9/16$ and $7/16$, using the
inequality $\tanh(\tanh x)\ge x/4$ on $[0,1]$ shown in the proof of
Lemma~\ref{lem:indexedgadget}.
\end{remark}

\subsubsection{A constant token gap from orthogonality}

\begin{lemma}[Dummy-key gap]
\label{lem:indexedgadget}
Given $m\ge1$ vectors $a_1,\ldots,a_m\in\{0,1\}^d$, $d\ge1$, define
\[
 L=\lceil\log_2(16m)\rceil,\qquad
 H=dL,\qquad n=H^2,\qquad r=dL^2.
\]
Repeat every coordinate of $-a_j$ exactly $r$ times to form $k_j$, and
repeat every coordinate of a query $b\in\{0,1\}^d$ exactly $r$ times
to form $q$.  Let $k_0=0$, $v_0=-1$, and $v_j=1$ for $j\ge1$.
Then
\[
 \langle q,k_j\rangle/\sqrt n=-L\langle a_j,b\rangle.
\]
If some $a_j$ is orthogonal to $b$, then
$\pi_{\mathcal C}(q)\ge1/2$.  If none is orthogonal, then
$\pi_{\mathcal C}(q)\le1/2-15/136<7/16$.
\end{lemma}

\paragraph{Proof idea.}
A zero dummy key gives a known reference weight.  An orthogonal vector matches that weight; without one, all real-key weights are exponentially suppressed.  Coordinate repetition realizes the needed penalty at standard scaling with sign entries.

\begin{proof}
There are $dr=d^2L^2=n$ coordinates, and $r/H=L$, proving the score
identity.  Put
\[
 W=\sum_{j=1}^m e^{-L\langle a_j,b\rangle}.
\]
The dummy has weight one, so the attention output is
$u=(W-1)/(W+1)$.  An orthogonal vector contributes one and gives
$W\ge1$, hence $u\ge0$.  Otherwise every dot product is at least one,
and
\[
 0\le W\le m e^{-L}\le m2^{-L}\le1/16.
\]
It follows that $u\le-15/17$.

For $0\le x\le1$, concavity of tanh on this interval and
$\tanh1\ge1/2$ give $\tanh x\ge x/2$.  The endpoint inequality follows
from $e^2\ge3$.  Applying this inequality twice gives
$\tanh(\tanh x)\ge x/4$.  Oddness and monotonicity therefore imply
\[
 \frac{1+\tanh(\tanh u)}2
 \le\frac12-\frac{15}{136}<\frac7{16}.
\]
All computations specifying the cache and query use integers.  The
reduction never evaluates an exponential, tanh, or an exact real number.
\end{proof}

\subsubsection{Absorbing every fixed polynomial preprocessing degree}

\paragraph{Proof idea.}
Partition the hard instance into small buckets so any fixed preprocessing degree is affordable.  Independent index copies amplify the token gap, and a global runtime cutoff converts the expected query guarantee into a bounded-time randomized decision algorithm.

\begin{proof}[Proof of Theorem~\ref{thm:indexedlower}]
Suppose the proposed data structure exists.  We may enlarge $a$ and
$b$ to integers and decrease $\epsilon$ to a positive rational if
necessary.  Set
\[
 \gamma=\frac1{2(a+1)},\qquad m=\lceil N^\gamma\rceil.
\]
Take a balanced randomized-OVH instance $A,B$, each of size $N$, in
dimension $d=\lceil c\log_2N\rceil$.  Partition $A$ into at most
$\lceil N/m\rceil$ buckets of size at most $m$.  Apply
Lemma~\ref{lem:indexedgadget} to each bucket.  Empty buckets are omitted.
For a smaller final bucket one may duplicate any of its vectors until
its size is $m$; orthogonality existence is unchanged.

Each index has $T=m+1$ entries and width
$n=(d\lceil\log_2(16m)\rceil)^2=O(\log^4N)=O(\log^4T)$.
The constant in this relation may depend on $c$ and $\gamma$.
For each $b\in B$, query every bucket.  There is an orthogonal pair
in $A\times B$ precisely when at least one bucket-query pair is in
the yes case of Lemma~\ref{lem:indexedgadget}.

Let
\[
 K=2048\left\lceil\log_2(24N^2)\right\rceil.
\]
Build $K$ mutually independent copies of each bucket index, and obtain
one sample for a fixed bucket-query pair from each independent copy.
For this fixed pair, the $K$ outputs are independent even when prefill
is randomized.  If the data structure can alter its index during a
query, record every changed cell and restore its old contents afterwards.
Every query to a given copy then starts in its original prefill state.
We bound the logging overhead after imposing a global time cutoff below;
this avoids any assumption on higher moments of a query's runtime.
No correctness assumption conditional on earlier outputs is being
introduced.

In the yes case each output has success probability at least $31/64$,
even with total variation error $1/64$.  In the no case it has success
probability at most $29/64$.  Classify a pair as yes if its empirical
mean is at least $15/32$.  Hoeffding's inequality bounds its error by
\[
 \exp(-2K(1/64)^2)\le\frac1{24N^2}.
\]
There are at most $N^2$ bucket-query pairs.  A union bound gives total
decision error at most $1/24$ before runtime truncation.  Outputs from
different pairs may be correlated; the union bound needs no
independence between pairs.

The expected preprocessing cost, including the independent copies, is
\[
 \widetilde O\left(\frac Nm m^a\right)
 =\widetilde O\left(N^{1+\gamma(a-1)}\right),
\]
and the total expected query cost is
\[
 \widetilde O\left(\frac{N^2}{m}m^{1-\epsilon}\right)
 =\widetilde O\left(N^{2-\gamma\epsilon}\right).
\]
Here the tilde hides only fixed powers of $\log N$.  Construction of
the explicit caches, formation of the explicit query vectors, sampling,
counting, and comparison are included.  In
particular the query vectors cost $O(n)$ bits each; this is absorbed
by the polylogarithmic factors.

Since $\gamma(a-1)<1/2$ and $\gamma\epsilon<1/4$, the total expected
original charged work is bounded by
$F(N)=O(N^{2-3\gamma\epsilon/4})$ for all sufficiently large $N$.
Simulate at most $24F(N)$ original charged steps, returning an
arbitrary answer on timeout.  Logging and undoing each memory write
then has at most an $O(\log(F(N)+2))=O(\log N)$ overhead in the
charged-address bit model: the log contains at most $24F(N)$ entries,
and the original address and changed bits were already charged.
The complete bounded simulation therefore takes
$O(F(N)\log N)=O(N^{2-\gamma\epsilon/2})$ bit operations.
Markov's inequality, applied to the original work, adds at most
$1/24$ error.  No bound on $\mathbb E[Q\log Q]$ is used.
We have obtained a
worst-case $O(N^{2-\gamma\epsilon/2})$-time randomized OV algorithm
with error at most $1/12$, contradicting randomized OVH for the
corresponding constant $c$.
\end{proof}

\subsubsection{Realization by row-norm-one transformer maps}

\begin{corollary}[The hard cache is realized by one legal block]
\label{cor:indexedblock}
The caches and queries in Theorem~\ref{thm:indexedlower} arise from a
one-block causal transformer with sign inputs, query, key, value, and
tanh row norms at most one, residual parameter $\alpha=1$, and
standard scaling $\beta=\sqrt n$.  The current position's own key and
value supply the dummy.  Consequently the lower bound applies to an
index over the already computed prefix keys and values.
\end{corollary}

\paragraph{Proof idea.}
Separate the raw coordinates used by cached keys from those used by the fresh query.  Shared half-sum rows select the appropriate role, and the current position contributes the same dummy key for every query.

\begin{proof}
Reserve two groups of $d$ input coordinates, a constant coordinate,
and a value marker.  There is room because
$n=d^2L^2\ge2d+2$.  Fill all unused input coordinates with $+1$.
At cache position $a_j$, put $2a_j-1$ in the first group, $-1$ in
the second, $+1$ in the constant coordinate, and $+1$ in the marker.
At the current query position $b$, put $-1$ in the first group,
$2b-1$ in the second, $+1$ in the constant coordinate, and $-1$ in
the marker.

For every one of the $r$ output rows assigned to original coordinate
$i$, let the query map be one half of the second-group coordinate
plus one half of the constant coordinate.  Let the key map be minus
one half of the first-group coordinate minus one half of the constant
coordinate.  These rows have $\ell_1$ norm one.  They produce exactly
the repeated query $b$ and repeated keys $-a_j$ in the lemma.
The current position has key zero.  Let the value map send the marker
to a chosen output coordinate and send zero to all other coordinates.
The selected attention output is $u$.  With $\alpha=1$, a unit tanh
row gives $\tanh u$, and the signed-coordinate head gives
\eqref{eq:indexedtarget}.  All coefficients lie in
$\{0,\pm1/2,\pm1\}$.

Only the first $m$ positions depend on the prefill input.  The key zero
and value $-1$ at the query position are the same for every query, so
an index may include this dummy during prefill.  No future query
information is used to build it.
\end{proof}

\subsubsection{A well-conditioned pre-norm additive realization}

\begin{corollary}[Normalization and additive residuals preserve the obstruction]
\label{cor:indexedprenorm}
The conclusion of Theorem~\ref{thm:indexedlower} also holds for a
single pre-norm additive attention block
\[
 h\longmapsto h+\operatorname{Att}(\operatorname{RMS}_0(h)),
\]
followed by a zero feed-forward update and a two-logit affine head
whose rows have $\ell_1$ norm at most one.  The RMS denominator before
attention equals one on every constructed input, and its squared
denominator after attention is at least $(n-1)/n$.  Thus the lower
bound does not rely on a vanishing normalization denominator.
\end{corollary}

\paragraph{Proof idea.}
All constructed raw streams are sign vectors, so normalization initially does nothing.  Store the attention output in the marker coordinate; an affine readout cancels the known residual offset and preserves a constant token gap.

\begin{proof}
Use the sign inputs and maps of Corollary~\ref{cor:indexedblock}.
Each sign vector has mean squared coordinate one, so
$\operatorname{RMS}_0(h)=h$.  Choose the marker itself as the output
coordinate of the value map.  The current marker is initially $-1$,
so after the additive attention update it is $-1+u$.  The constant
coordinate stays equal to one.  Their quarter-sum is
$z=u/4\in[-1/4,1/4]$.  Use logits $\pm z$; the token probability is
$(1+\tanh(u/4))/2$.  This dyadic readout also satisfies the bounded-logit
promise on the entire sign cube: the updated marker has magnitude at
most two and the unchanged constant coordinate has magnitude one,
so the quarter-sum has magnitude at most $3/4$ even outside the
restricted hard family.

For $0\le x\le1/4$, one has
$\tanh x\ge x-x^3/3\ge(7/8)x$.
Indeed $0\le\tanh x\le x$ gives
$\tanh'(x)=1-\tanh^2x\ge1-x^2$, and integration proves the first
inequality.  The second follows from $x^2\le1/16$.
The no case has $u\le-15/17$, hence
\[
 (1+\tanh(u/4))/2\le1/2-105/1088<7/16;
\]
the yes case is at least $1/2$.  The same amplification proof
therefore applies.  All other
coordinates are unchanged signs.  At every position after attention,
their contribution to the mean square is $(n-1)/n$.  Set the
subsequent tanh feed-forward map to zero.  Its normalization is
well-defined with the displayed floor, and the update is zero.
\end{proof}

For a fixed positive dyadic stabilizer $\epsilon_0$, the same construction
rescales query and key vectors by $1/\sqrt{1+\epsilon_0}$.
Replacing $L$ by an integer at least
$(1+\epsilon_0)\lceil\log_2(16m)\rceil$ restores the score gap.
The value amplitude is reduced by the fixed factor
$1/\sqrt{1+\epsilon_0}$, leaving a positive constant token gap.
The number of independent copies changes by a constant depending on
$\epsilon_0$.  All matrix coefficients remain rational constants.
The reduction does no normalization arithmetic; the exact sampler
performs it.
For exact sampling this proves the conclusion for every fixed
$\epsilon_0>0$.  A fixed total-variation tolerance must be chosen
smaller than a constant multiple of
$(1+\epsilon_0)^{-1/2}$; the particular tolerance $1/64$ continues
to work, with the displayed thresholds, whenever $0\le\epsilon_0\le1$.

\subsubsection{\texorpdfstring{Width $O(\log^2 T)$}{Quadratic logarithmic width} from Hamming-distance gaps}

The elementary orthogonality construction used a penalty for a single
violated coordinate.  A known gap reduction gives a stronger dimension
bound.  We state the imported hardness fact precisely before using it.

\begin{lemma}[Gap Hamming hardness, extracted from Rubinstein]
\label{lem:gaphamminghardness}
Assume randomized OVH.  For every $\rho>0$, there are fixed positive
constants $\zeta,c_0,C_0$ such that the following promise problem has
no randomized $O(N^{2-\rho})$-time algorithm with error at most $1/3$.
Given two lists $A,B$ of $N$ vectors in $\{0,1\}^d$, with
$c_0\log_2N\le d\le C_0\log_2N$, and an integer
$d/4\le t\le d/2$, distinguish
\[
 \min_{a\in A,b\in B}\operatorname{dist}_H(a,b)\le t
 \quad\hbox{from}\quad
 \min_{a\in A,b\in B}\operatorname{dist}_H(a,b)\ge(1+\zeta)t.
\]
The gap constant may be chosen rational.
\end{lemma}

This is the Hamming part of Rubinstein's deterministic reduction in
the proof of Theorem~4.1 of~\cite{rubinstein2018ann}.  In its notation,
the output dimension is $d=2T'm_0$, the yes distance is
$t=(T'-1)m_0$, and the no distance is at least $T'm_0$, where $T'\ge2$
is constant.  These formulas give $d/4\le t<d/2$ and a fixed rational
gap.  The verifier's random tapes become vector coordinates; they are
enumerated by the reduction.  Its deterministic construction therefore
also transfers hardness against randomized algorithms.  Duplicating
vectors balances the two output sets without changing the distances.

\begin{theorem}[Indexed lower bound at quadratic logarithmic width]
\label{thm:indexedlowerlogtwo}
Under randomized OVH, Theorem~\ref{thm:indexedlower} remains true with
$n\le C\log^2(T+2)$, allowing $q,k_i\in\{-1,0,1\}^n$.
The queries and all non-dummy keys can in fact have every coordinate
in $\{-1,1\}$, so their Euclidean norms equal $\sqrt n$.
The total-variation tolerance $1/64$ is unchanged.
\end{theorem}

\paragraph{Proof idea.}
A Hamming-distance gap turns the required score penalty into a constant number of coordinate repetitions.  The same dummy-key separation and bucket amplification then give the sharper width bound.

\begin{proof}
Fix proposed preprocessing degree $a$ and query saving $\epsilon$,
and set $\gamma=1/[2(a+1)]$ as before.  Use
Lemma~\ref{lem:gaphamminghardness} with
$\rho=\gamma\epsilon/2$; let its constants be $\zeta,c_0,C_0$.
Take one of its instances of size $N$, and partition $A$ into buckets
of size $m=\lceil N^\gamma\rceil$, duplicating the final bucket's
elements if necessary.

For the known threshold $t$, fix a vector $c_t\in\{-1,1\}^d$
with exactly $t$ positive entries.  Given $a,b\in\{0,1\}^d$, form
\[
 \widehat k_a=(2a-1,c_t),\qquad
 \widehat q_b=(2b-1,\boldsymbol1)
 \quad\hbox{in }\{-1,1\}^{2d}.
\]
Their dot product is exactly
\begin{equation}
 \langle\widehat q_b,\widehat k_a\rangle
 =d-2\operatorname{dist}_H(a,b)+(2t-d)
 =2(t-\operatorname{dist}_H(a,b)).
 \label{eq:gaphammingdot}
\end{equation}
Let $L=\lceil\log_2(16m)\rceil$ and choose
\[
 H=\max\{1,\lceil L/(2\zeta t)\rceil\},\qquad
 r=2dH^2,\qquad n=(2dH)^2.
\]
The gap constant is a fixed rational number, so this choice uses only
rational arithmetic and an integer ceiling.  Repeat every coordinate
of $\widehat k_a$ and $\widehat q_b$ exactly $r$ times.  Since
$2dr=n$ and $r/\sqrt n=H$, the standard-scaled score is
\[
 \frac{\langle q_b,k_a\rangle}{\sqrt n}
 =2H(t-\operatorname{dist}_H(a,b)).
\]
Append a zero dummy key with value $-1$ and give all real keys value
$+1$.

In a yes bucket-query pair some real key has score at least zero,
so its total real-key weight $W$ is at least one.  In a no instance,
every real score is at most $-2H\zeta t\le-L$, and every bucket has
$W\le m e^{-L}\le1/16$.  The dummy-key argument gives exactly the
same constant token gap as Lemma~\ref{lem:indexedgadget}.

Here $t\ge d/4\ge(c_0/4)\log_2N$, whereas
$L\le\gamma\log_2N+O(1)$.  Therefore $H$ is bounded by a constant
depending only on $\gamma,\zeta,c_0$.  Consequently
\[
 n=O(d^2)=O(\log^2N)=O(\log^2T).
\]
All input coordinates are $0$ or signs; the scaling denominator
$\sqrt n=2dH$ is an integer.

Use the independent prefill copies, sample threshold, union bound,
and runtime truncation from Theorem~\ref{thm:indexedlower}.  A yes
instance has at least one yes bucket-query pair.  A no instance has
only no pairs.  Under the promise, these are the only cases needed;
intermediate distances in a yes instance can produce arbitrary
additional yes classifications without affecting the decision.
The expected time is
\[
 \widetilde O\left(N^{1+\gamma(a-1)}
                 +N^{2-\gamma\epsilon}\right),
\]
which after truncation is
$O(N^{2-\gamma\epsilon/2})$.  This contradicts
Lemma~\ref{lem:gaphamminghardness}.
\end{proof}

This matches the exponent of the straightforward scan.  At width
$n=O(\log^2T)$ and input bit lengths polylogarithmic in $T$,
Lemma~\ref{lem:newcachescan} samples \eqref{eq:indexedtarget} exactly
in $T\operatorname{polylog}T$ expected bit operations.  Thus the lower
and upper bounds agree in the power of $T$, under the stated
hypothesis.

\begin{corollary}[Quadratic-logarithmic hard blocks]
\label{cor:indexedlogtwoblock}
The row-norm-one block realization and the well-conditioned pre-norm
additive realization in Corollaries~\ref{cor:indexedblock} and
\ref{cor:indexedprenorm} also hold for the
$n=O(\log^2T)$ lower bound of Theorem~\ref{thm:indexedlowerlogtwo}.
\end{corollary}

\paragraph{Proof idea.}
Opposite sign pairs encode a coordinate at its intended position and give zero at the other role.  Fixed padding and the marker coordinate complete the Hamming-gap gadget using only rows of absolute sum one.

\begin{proof}
For $d\ge2$, reserve $4d+2$ source coordinates, which fit in
$n=(2dH)^2$.  For each $i\le d$, use two key-source coordinates and
two query-source coordinates.  At a cache position $a$, the key pair
is $(2a_i-1,1-2a_i)$ and the query pair is $(1,1)$.  At the current
position $b$, the key pair is $(1,1)$ and the query pair is
$(2b_i-1,1-2b_i)$.  Half the difference of each pair gives the
required bipolar key or query coordinate; it is zero on the other
role.  Reserve a constant coordinate $c=1$ and a marker $v$ that is
$+1$ on cache positions and $-1$ on the query position.

The $i$th fixed padding coordinate of a key is
$(c_t)_i(c+v)/2$, and the corresponding query padding coordinate is
$(c-v)/2$.  Thus cache keys have padding $c_t$, the query has padding
$\boldsymbol1$, and its own key is zero.  Replicate every one of
these $2d$ output rows $r$ times.  Each row has absolute sum one.
The marker supplies the scalar value.  All raw source coordinates
are signs, including the unused coordinates, which are set to one.

The remainder of Corollary~\ref{cor:indexedblock}, or respectively
Corollary~\ref{cor:indexedprenorm}, now applies verbatim.  The
asymptotic hardness instances have $d\ge2$; finitely many smaller
instances have no effect on the contradiction.
\end{proof}

